% Options for packages loaded elsewhere
\PassOptionsToPackage{unicode}{hyperref}
\PassOptionsToPackage{hyphens}{url}
%
\documentclass[
]{article}
\usepackage{amsmath,amssymb}
\usepackage{iftex}
\ifPDFTeX
  \usepackage[T1]{fontenc}
  \usepackage[utf8]{inputenc}
  \usepackage{textcomp} % provide euro and other symbols
\else % if luatex or xetex
  \usepackage{unicode-math} % this also loads fontspec
  \defaultfontfeatures{Scale=MatchLowercase}
  \defaultfontfeatures[\rmfamily]{Ligatures=TeX,Scale=1}
\fi
\usepackage{lmodern}
\ifPDFTeX\else
  % xetex/luatex font selection
\fi
% Use upquote if available, for straight quotes in verbatim environments
\IfFileExists{upquote.sty}{\usepackage{upquote}}{}
\IfFileExists{microtype.sty}{% use microtype if available
  \usepackage[]{microtype}
  \UseMicrotypeSet[protrusion]{basicmath} % disable protrusion for tt fonts
}{}
\makeatletter
\@ifundefined{KOMAClassName}{% if non-KOMA class
  \IfFileExists{parskip.sty}{%
    \usepackage{parskip}
  }{% else
    \setlength{\parindent}{0pt}
    \setlength{\parskip}{6pt plus 2pt minus 1pt}}
}{% if KOMA class
  \KOMAoptions{parskip=half}}
\makeatother
\usepackage{xcolor}
\usepackage[margin=1in]{geometry}
\usepackage{graphicx}
\makeatletter
\def\maxwidth{\ifdim\Gin@nat@width>\linewidth\linewidth\else\Gin@nat@width\fi}
\def\maxheight{\ifdim\Gin@nat@height>\textheight\textheight\else\Gin@nat@height\fi}
\makeatother
% Scale images if necessary, so that they will not overflow the page
% margins by default, and it is still possible to overwrite the defaults
% using explicit options in \includegraphics[width, height, ...]{}
\setkeys{Gin}{width=\maxwidth,height=\maxheight,keepaspectratio}
% Set default figure placement to htbp
\makeatletter
\def\fps@figure{htbp}
\makeatother
\setlength{\emergencystretch}{3em} % prevent overfull lines
\providecommand{\tightlist}{%
  \setlength{\itemsep}{0pt}\setlength{\parskip}{0pt}}
\setcounter{secnumdepth}{-\maxdimen} % remove section numbering
\ifLuaTeX
  \usepackage{selnolig}  % disable illegal ligatures
\fi
\usepackage{bookmark}
\IfFileExists{xurl.sty}{\usepackage{xurl}}{} % add URL line breaks if available
\urlstyle{same}
\hypersetup{
  pdftitle={README},
  hidelinks,
  pdfcreator={LaTeX via pandoc}}

\title{README}
\author{}
\date{\vspace{-2.5em}}

\begin{document}
\maketitle

\textbf{\emph{Data Folders}} 1. \textbf{LICOR} This folder includes all
of the raw and standardized data from infrared gas analysis lab
experiments. This data is further organized by crust type, including
moss cultures, combined CO2/water curves on cultures, combined CO2/water
curves on cyanobacteria crust, cyanobacteria crust single factor curves,
and lichen crust curves. Within a given folder, there are raw data files
outputted directly from the LICOR-6800. Defined below are the relevant
columns:\\
- elapsed (seconds)\\
- A\_scaled (mumol/m2/s): carbon assimilation scaled by dry weight and
mass of sample\\
- A (mumol/kg/s): carbon assimilation recorded by LICOR machine\\
- Qin (mumol/m2/s): light level observed by LICOR machine\\
- CO2\_r (ppm): CO2 level observed by LICOR machine\\
- Tair (degC): temperature observed by LICOR machine\\

These files are used as a foundation to develop relevant mathematical
relationships. This data has been cleaned and scaled appropriately by
dry weights and masses. These data were collected by Asha Kannan
(2025-2026), and cleaned with protocol developed by Kirsten Coe.\\

\begin{enumerate}
\def\labelenumi{\arabic{enumi}.}
\setcounter{enumi}{1}
\tightlist
\item
  \textbf{NMEG} This folder includes all the flux tower data files, with
  both gaps and gapfilled, available. Our model pulls the gapfilled
  files. All files are named with the following format:
  US-Mpj\_{[}year{]}\_gapfilled.txt. Defined below are the relevant
  columns:\\
\end{enumerate}

\begin{itemize}
\tightlist
\item
  TIMESTAMP\_START (YYYYMMDDHHMM) and TIMESTAMP\_END (YYYYMMDDHHMM):
  used as timesteps in the model\\
\item
  PPFD\_IN (mumol/m2/s): light used in the model\\
\item
  GPP\_DT (mumol/m2/s): ecosystem GPP (gapfilled) in model\\
\item
  RECO\_DT (mumol/m2/s): ecosystem respiration (gapfilled) in model\\
\item
  TA\_F (degC): temperature recorded by flux tower\\
\item
  P\_F\_FLAG (--)\\
\end{itemize}

These files are used in DrylandBSCModel\_June2026.R in additional to
angular\_interpolation\_building.R. These data were provided by Rae
DeVan \& Marcy Litvak at the University of New Mexico.\\

\begin{enumerate}
\def\labelenumi{\arabic{enumi}.}
\setcounter{enumi}{2}
\tightlist
\item
  \textbf{monthly\_footprint\_climatology\_weight\_map}\\
  This folder includes all the monthly footprint climatology .tif files
  downloaded from Chu et al.~(2021) Supplementary Dataset. See more
  information here: \url{https://zenodo.org/records/4015350}. All files
  are in the following format:
  US-Mpj\_{[}YYYY{]}\emph{{[}MM{]}}{[}DAY/NIGHT{]}\_fpt\_weight.tif.\\
\end{enumerate}

These files are used in angular\_interpolation\_building.R.\\

\begin{enumerate}
\def\labelenumi{\arabic{enumi}.}
\setcounter{enumi}{3}
\tightlist
\item
  \textbf{PJ\_soil\_allyears\_30min.csv} This file contains relevant
  soil moisture and soil temperature information from sensors at the
  US-Mpj site. Columns are defined in this format for temperature
  sensors: SOILT\_{[}P,O,J{]}{[}1,2,3{]}\_{[}5,10,30{]}\_AVG, where
  P/O/J relates to the relevant canopy cover (pinyon pine, open, or
  juniper), 1/2/3 refers to the unit number, and 5/10/30 refers to the
  depth of the relevant sensor, either 5, 10, or 30 cm below the
  surface.
\end{enumerate}

This file was used as environmental inputs in
DrylandBSCModel\_June2026.R. These data were provided by Rae DeVan \&
Marcy Litvak at the University of New Mexico.

\begin{enumerate}
\def\labelenumi{\arabic{enumi}.}
\setcounter{enumi}{4}
\item
  \textbf{mpj\_transects.xslx} This file contains vegetation
  identification as measured through the LPI methodology in the field by
  Kirsten Coe and Asha Kannan in October 2025. These data are used to
  build the vegetation raster in angular\_interpolation\_building.R.
\item
  \textbf{results\_coverfrac.csv} This file contains the results of the
  cover fractions of different vegetation types that is an output of
  angular\_interpolation\_building.R. Defined below are the relevant
  columns:
\end{enumerate}

\begin{itemize}
\tightlist
\item
  year, month, period (day or night), w\_cyano (weight of cyanobacteria
  crust), w\_lich (weight of lichen crust), w\_moss (weight of moss
  crust), w\_midveg (weight of mid-level vegetation), w\_subcanopy
  (weight of pinyon/juniper but not as a canopy), w\_canopy (weight of
  overhanging canopy), w\_abiotic (weight of nonphotosynthetic material,
  eg hardware, fallen leaves, etc)\\
\end{itemize}

\textbf{\emph{Data Analysis}}\\
1. \textbf{RollingSlope.R}\\
Authors: Casey Schine and Lauren Marsala\\
Purpose: time-series method that identifies points where the calculated
slope changes too drastically compared to other nearby data points. Used
to clean outliers in LICOR data. Arguments: x, y, width\\

\begin{enumerate}
\def\labelenumi{\arabic{enumi}.}
\setcounter{enumi}{1}
\tightlist
\item
  \textbf{PhotosynthesisFunctions.R}\\
  Authors: Casey Schine and Asha Kannan, adapted from Webb et
  al.~(1974); Farquhar et al.~(2001); and Ye et al.~(2021)\\
  Purpose: Functions for temperature, light, and water developed to
  describe biocrust photosynthesis response\\
\end{enumerate}

Temperature: calc\_Pg\_fxn\_T\_Tk\_rowwise\\
Arguments: temp (measured soil temperature); T\_k (controls rate of
rise); Pgmax (max photosynthesis)\\
Output: out\_P\\
Temperature: f\_H\\
Arguments: Tv (measured temperature); delta\_S (entropy); delta\_Hd
(deactivation enthalpy); R\_gas (universal gas constant)\\
Output: f\_H (high-temperature deactivation term)\\
Dark respiration: calc\_RD\_fxn\_T\_rowwise\\
Arguments: temp (measured temperature); r\_r (temperature sensitivity
coefficient); R\_o (base respiration at 0C)\\
Output: out\_R\\

Light: calc\_PN\_fxn\_I\_rowwise\\
Arguments: PAR (measured light level); f\_Io\_Icomp (apparent quantum
yield); b (photoinhibition coefficient); g (curvature coefficient);
Icomp (light compensation point)\\
Output: PN\\

Water: low\_water\_response\_fxn\\
Arguments: WC (normalized water content); WC\_k (threshold water
content)\\
Output: low-water activation (out\_P)\\
Water: high\_water\_response\_fxn\\
Arguments: WC (normalized water content); toggled intercepts for
high-water activation\\
Output: high-water activation (out\_P)\\

\begin{enumerate}
\def\labelenumi{\arabic{enumi}.}
\setcounter{enumi}{2}
\item
  \textbf{ye\_et\_al\_light\_fit.xlsx}\\
  This solver spreadsheet, developed by Casey Schine, was adapted from
  Ye et al.~(2021) for fitting light curves.\\
  Inputs: Bin LICOR data for each light level, then average per bin such
  that there is one value for carbon assimilation per light level.
  Instructions are in Sheet 2 of the spreadsheet.\\
  Outputs: Save all parameters that are outputted by the spreadsheet,
  including the sum of squares.\\
\item
  \textbf{angular\_interpolation\_building.R}\\
  Script developed by Asha Kannan, to build the vegetation raster based
  on field-collected spatial vegetation cover data.\\
  Packages required: ``readxl'', ``dplyr'', ``tidyr'', ``ggplot2'',
  ``terra'', ``stringr'', ``av'', ``here''\\
  Inputs required: mpj\_transects.xlsx;
  monthly\_footprint\_climatology\_weight\_map\\
  Output: results\_coverfrac.csv\\
\item
  \textbf{DrylandBSCModel\_June2026.R}\\
  Script developed by Asha Kannan, as ecophysiology-based model
  framework.\\
  Packages required: ``tidyverse'', ``readxl'', ``minpack.lm'',
  ``lubridate'', ``zoo'', ``patchwork''\\
  Inputs required: NMEG, PJ\_soil\_allyears\_30min.csv,
  results\_coverfrac.csv\\
  Outputs: model\_daily, model\_monthly, model\_annual\\
\end{enumerate}

\end{document}
